\documentclass[11pt,a4paper]{article}
\usepackage{notes}
\usepackage{fancyhdr}

\pagestyle{fancy}
\fancyhf{}
\lhead{Geophysics 3A: Potential Fields \& Geothermics}
\rhead{Week 7 Workshop}
\cfoot{\thepage}

\title{Week 7 -- Workshop Plan\\Fourier Transforms and Potential Fields}
\author{Geophysics 3A: Potential Fields \& Geothermics}
\date{\today}

\begin{document}
\maketitle

\section*{Workshop Overview}

This workshop develops intuition for Fourier methods in potential-field geophysics
and applies them to gravity and magnetic interpretation.
The emphasis is on understanding how sampling, filtering, and phase
control what geological information can be recovered.

\section*{Learning Objectives}

By the end of this workshop, students should be able to:
\begin{itemize}
  \item Relate spatial wavelength to source depth and geometry
  \item Recognize aliasing and understand its consequences
  \item Interpret gravity and magnetic spectra in geological terms
  \item Apply upward and downward continuation appropriately
  \item Explain the purpose and limitations of reduction to the pole (RTP)
\end{itemize}

\section*{Timeline and Activities}

\subsection*{0--15 min \quad Introduction and framing}
\begin{itemize}
  \item Why Fourier methods appear everywhere in potential-field geophysics
  \item Review of key ideas: wavelength, sampling, phase
  \item Overview of the day's activities
\end{itemize}

\subsection*{15--40 min \quad Demo 1: What a Fourier transform does}
\begin{itemize}
  \item Decomposing a signal into frequency, amplitude, and phase
  \item Phase as a measure of shift and geometry
  \item Padding and windowing
\end{itemize}

\noindent\textbf{Key emphasis:} Fourier transforms reorganize information; they do not interpret it.

\subsection*{40--75 min \quad Demo 2: Spatial sampling and aliasing}
\begin{itemize}
  \item Synthetic fold and Basin-and-Range gravity models
  \item Station spacing and Nyquist wavelength
  \item Aliasing as coherent but incorrect structure
\end{itemize}

\noindent\textbf{Key emphasis:} Sampling decisions set a hard limit on geological resolution.

\subsection*{75--95 min \quad Real data case study: Nevada gravity}
\begin{itemize}
  \item Complete Bouguer gravity of Nevada
  \item E--W and N--S power spectra
  \item Identification of Basin-and-Range wavelengths
\end{itemize}

\subsection*{95--115 min \quad Isostatic correction and filtering}
\begin{itemize}
  \item Airy isostasy as a physical low-pass filter
  \item Isostatic residual gravity
  \item High-pass filtering and wavelength selection
\end{itemize}

\noindent\textbf{Key emphasis:} Filtering encodes geological assumptions.

\subsection*{115--135 min \quad Inversion example}
\begin{itemize}
  \item Parker--Oldenburg sediment thickness inversion
  \item Interpretation and limitations
  \item Effect of aliasing on inversion results
\end{itemize}

\subsection*{135--165 min \quad Magnetic interpretation: RTP and maar geometry}
\begin{itemize}
  \item Magnetic anomalies as vector fields
  \item Cone/maar geometry and asymmetric anomalies
  \item Reduction to the pole as a phase operation
\end{itemize}

\noindent\textbf{Key emphasis:} RTP rotates phase; it does not filter noise.

\subsection*{165--180 min \quad Synthesis and discussion}
\begin{itemize}
  \item Gravity vs magnetics: robustness and fragility
  \item When Fourier methods help—and when they mislead
  \item Connection to upcoming assignments and field interpretation
\end{itemize}

\section*{Key Takeaways}

\begin{itemize}
  \item Fourier methods are powerful but unforgiving
  \item Sampling precedes interpretation
  \item Phase carries geometry
  \item Geological assumptions are embedded in every filter
\end{itemize}

\end{document}